\section*{Related work}
\subsection*{Transcriptomic studies of liver fibrosis}
Hepatic transcriptomic studies have identified molecular signatures associated with steatohepatitis and fibrosis severity. Hoang et al. analysed 78 human liver biopsies spanning histological severity and identified progressively regulated genes and heterogeneous disease-associated expression profiles \cite{ref4}. Their study provides a direct precedent for investigating expression in relation to ordered fibrosis stages. Govaere et al. characterised transcription across the NAFLD spectrum in a discovery cohort of 206 patients with NAFLD, identified a 25-gene signature associated with fibrosing steatohepatitis, and evaluated it in an independent replication cohort. Serum-protein measurements and functional experiments provided additional evidence for selected findings \cite{govaere2020}. These studies establish that stage-associated expression signatures and molecular heterogeneity are existing areas of investigation. They also distinguish discovery of an expression association from clinical or functional validation of a candidate.

\subsection*{Cross-cohort integration of liver transcriptomic data}
Piras and DiStefano integrated ten hepatic RNA-sequencing and microarray datasets comprising 1,058 samples using a random-effects meta-analysis. Comparisons among normal liver, MASLD and MASH were complemented by genetic-association and coexpression-network analyses \cite{piras2024}. Their assessment of coexpression-module preservation supported network reproducibility, while experimental validation of the proposed regulatory roles remained a separate requirement. This work demonstrates the value of integrating cohorts and organising associations into biological modules.

Maiers et al. analysed hepatic transcriptomic data from approximately 1,000 patients with MASLD and identified 846 differentially regulated genes in an advanced-versus-mild fibrosis comparison, F3--F4 versus F0--F1. Cell-type deconvolution and ligand--receptor analyses provided additional biological context \cite{maiers2026}. Their principal grouped-stage comparison addresses conserved features of advanced fibrosis. Retaining F0--F4 separately provides a complementary perspective on intermediate-stage patterns and adjacent-stage differences, rather than establishing superiority over the grouped comparison.

\subsection*{Fibrosis-stage expression patterns and molecular heterogeneity}
Previous studies have organised heterogeneity through disease-associated expression profiles, patient clustering, coexpression modules and inferred cellular composition \cite{ref4,govaere2020,piras2024,maiers2026}. These approaches address different levels of biological organisation and should not be treated as interchangeable. Patient clustering identifies groups of samples, whereas clustering genes by their stage-expression profiles identifies sets of genes with similar patterns across ordered stages.

The present analysis builds on this literature by grouping fibrosis-associated genes according to their five-stage profiles and examining adjacent-stage contrasts alongside those patterns. Because the measurements are predominantly cross-sectional, an early or late expression pattern describes differences among patients at different stages rather than an observed molecular sequence within an individual. Gene assignments and average patterns also require separate stability assessments when studies or patients are resampled.

\subsection*{Knowledge graphs for biological interpretation}
Biomedical knowledge graphs connect molecular findings with heterogeneous biological information through explicit entity and relationship types. Himmelstein et al. developed Hetionet by integrating 29 public resources, demonstrating how genes, diseases, pathways and other entities can be represented within a common network \cite{himmelstein2017}. OptimusKG extends this resource-oriented approach through a labelled property graph assembled from structured and semi-structured biomedical sources. Its source resource retains entity-specific properties, cross-references and provenance across molecular, anatomical, clinical and environmental domains; the resource is described in a preprint \cite{ref5}.

These resources provide a foundation for contextualising fibrosis-associated genes through curated pathway and biological-process relationships. The properties retained in a study-specific graph depend on its import procedures and must be distinguished from the complete source resource. Likewise, annotation sharing and graph connectivity support interpretation but do not independently establish causal regulation or validate candidate genes.

\subsection*{Research gap and study rationale}
Previous studies have characterised fibrosis-associated expression signatures, molecular heterogeneity, cross-cohort reproducibility and biological networks. Building on these contributions, the present study examines how cross-study association evidence, expression profiles across F0--F4, adjacent-stage differences and curated biological annotations can be organised within a common framework for candidate-gene prioritisation. Its contribution lies in integrating these complementary evidence types while reporting their uncertainty and internal robustness.

The framework retains separate fibrosis-stage and ordered disease-group analyses, combines explicitly defined prioritisation features, and links prioritised genes to curated biological context. It supports investigation of fibrosis-associated transcriptional organisation. Independent longitudinal and functional studies remain necessary to establish temporal progression and causal biological roles, while clinical prediction would require a separately designed and independently evaluated model.

